\documentclass[11pt,letterpaper]{article}
\usepackage[T1]{fontenc}
\usepackage[utf8]{inputenc}
\usepackage{lmodern}
\usepackage[margin=1in,headheight=15pt]{geometry}
\usepackage{microtype}
\usepackage{amsmath,amssymb,amsthm,mathtools}
\usepackage{booktabs,array,longtable,aliascnt}
\usepackage{enumitem,fancyhdr,xcolor,xurl,hyperref}
\hypersetup{colorlinks=true,linkcolor=blue!40!black,citecolor=blue!40!black,
 urlcolor=blue!40!black,pdftitle={Exponential Relations over Omnific Integers},
 pdfauthor={AI-assisted research manuscript},
 pdfsubject={Exact toric kernels, decidable fragments, and elementary substructures}}
\usepackage[nameinlink,noabbrev]{cleveref}
\numberwithin{equation}{section}
\newtheorem{theorem}{Theorem}[section]
\newaliascnt{lemma}{theorem}
\newtheorem{lemma}[lemma]{Lemma}
\aliascntresetthe{lemma}
\newaliascnt{proposition}{theorem}
\newtheorem{proposition}[proposition]{Proposition}
\aliascntresetthe{proposition}
\newaliascnt{corollary}{theorem}
\newtheorem{corollary}[corollary]{Corollary}
\aliascntresetthe{corollary}
\theoremstyle{definition}
\newaliascnt{definition}{theorem}
\newtheorem{definition}[definition]{Definition}
\aliascntresetthe{definition}
\newaliascnt{example}{theorem}
\newtheorem{example}[example]{Example}
\aliascntresetthe{example}
\newtheorem{question}{Research question}
\theoremstyle{remark}
\newaliascnt{remark}{theorem}
\newtheorem{remark}[remark]{Remark}
\aliascntresetthe{remark}
\crefname{lemma}{lemma}{lemmas}
\crefname{proposition}{proposition}{propositions}
\crefname{corollary}{corollary}{corollaries}
\crefname{definition}{definition}{definitions}
\crefname{example}{example}{examples}
\crefname{remark}{remark}{remarks}
\crefname{question}{research question}{research questions}
\newcommand{\No}{\mathbf{No}}
\newcommand{\Oz}{\mathbf{Oz}}
\newcommand{\SC}{\mathbf{No}[i]}
\newcommand{\Z}{\mathbb Z}
\newcommand{\Q}{\mathbb Q}
\newcommand{\R}{\mathbb R}
\newcommand{\C}{\mathbb C}
\newcommand{\N}{\mathbb N}
\newcommand{\A}{\mathbb A}
\newcommand{\Qb}{\overline{\mathbb Q}}
\newcommand{\J}{\mathbb J}
\newcommand{\M}{\mathfrak M}
\newcommand{\Ealg}{\mathcal E_{\mathrm{alg}}}
\newcommand{\Lrat}{\mathcal L_{\mathrm{rat}}}
\newcommand{\Lalg}{\mathcal L_{\mathrm{alg}}}
\newcommand{\LPr}{\mathcal L_{\mathrm{Pr}}}
\newcommand{\Lvec}{\mathcal L_{\A\mathrm{-vec}}}
\newcommand{\ct}{\operatorname{ct}}
\newcommand{\pp}{\operatorname{pp}}
\newcommand{\supp}{\operatorname{supp}}
\newcommand{\Exp}{\operatorname{Exp}}
\newcommand{\Frac}{\operatorname{Frac}}
\newcommand{\trdeg}{\operatorname{trdeg}}
\newcommand{\spanQ}{\operatorname{span}_{\Q}}
\newcommand{\Aut}{\operatorname{Aut}}
\newcommand{\Th}{\operatorname{Th}}
\newcommand{\bal}{\operatorname{Bal}}
\newcommand{\Pure}{\mathsf{Pure}}
\newcommand{\Int}{\mathsf{Int}}
\newcommand{\ctgraph}{\mathsf{CT}}
\newcommand{\repoSHA}{\nolinkurl{958b5c4865819bd55ea1f5ffc050282aea7ef570}}
\setlist[enumerate]{leftmargin=2em,itemsep=3pt,topsep=5pt}
\setlength{\parskip}{3pt}
\setlength{\parindent}{1.2em}
\emergencystretch=2em
\pagestyle{fancy}
\fancyhf{}
\fancyhead[L]{\small Exponential relations over omnific integers}
\fancyhead[R]{\small Research manuscript}
\fancyfoot[C]{\thepage}
\renewcommand{\headrulewidth}{0.3pt}
\setcounter{tocdepth}{2}
\begin{document}
\begin{titlepage}
\centering
\vspace*{0.35in}
{\large\scshape Surreal algebra and model theory\par}
\vspace{0.3in}
{\LARGE\bfseries Exponential Relations\par over Omnific Integers\par}
\vspace{0.24in}
{\large Exact toric kernels, decidable fragments,\par and one-scale elementary cores\par}
\vspace{0.32in}
{23 September 2026\par}
\vspace{0.35in}
\begin{minipage}{0.92\textwidth}
\begin{abstract}
Let $\J$ be the purely infinite surreals, so that Conway's omnific integers are
$\Oz=\Z\oplus\J$. On the additive domain $\Ealg=\Qb+\J\subseteq\SC$, define
$\Exp(a+p)=e^a\exp(p)$. Combining Gonshor's monomial theorem with classical
Lindemann--Weierstrass gives an injective exponential group-algebra map. We derive
an exact description of all Laurent-polynomial relations among finitely many
exponentials, their transcendence degree, and the ordinary complex constants in
the field they generate.

The principal model-theoretic application concerns omnific-integer variables and
finite exponential sums with algebraic coefficients, algebraic shifts, and real
algebraic linear slopes. Their zero relations admit an effective balanced-partition
elimination. With addition, order, and congruences, the resulting first-order
language is decidable. Its rational-slope reduct is a definitional expansion of
Presburger arithmetic. The algebraic-slope structure has a definable decomposition
into ordinary Presburger integers and a nonzero ordered vector space over the real
algebraic numbers. We classify all its elementary substructures: they are exactly
$\Z+W$ for nonzero algebraic-vector subspaces $W\subseteq\J$. In particular,
$\Z+\A u$ is a countable elementary core for every positive $u\in\J$.

We also prove a rectangular normal form for definable sets, finite-image results
for maps between the two components, and an undecidability boundary obtained by
adjoining the single nonlinear relation $\exp(xy)=\exp(z)$. Additional results
classify certain exponential equations with arbitrary positive real bases by
ordinary integer fibers and purely infinite linear families. With just one
irrational scalar-graph predicate, we prove that the theory has exactly the Turing
degree of the rational cut of its slope.
\end{abstract}
\end{minipage}
\vfill
\begin{minipage}{0.92\textwidth}\small
\textbf{Status and scope.} An AI-assisted, unrefereed research manuscript. Complete
arguments are supplied relative to explicitly stated classical inputs. The main
classification, computability, and decidability results are proposed contributions; priority is
not certified. No named classical conjecture, full surreal exponential theory,
or unrestricted surcomplex exponential problem is claimed solved. Finite
verification code is included, but no Lean verification is claimed.
\end{minipage}
\end{titlepage}
\tableofcontents
\clearpage
\section{Introduction and principal results}\label{sec:intro}
\subsection{A restricted language with an exact answer}
An omnific integer may contain an arbitrarily complicated set-sized normal form,
yet some finite exponential questions about omnific integers see only a surprisingly
small part of that complexity. The point of this article is to identify a language
in which this observation becomes an exact structural theorem, rather than a
heuristic about dominant terms.

Write $\A=\Qb\cap\R$ for the ordered field of real algebraic numbers. Every
$x\in\Oz$ has a unique expression
\begin{equation}\label{eq:introsplit}
 x=n+p,\qquad n\in\Z,\quad p\in\J.
\end{equation}
Here ``purely infinite'' means that every exponent in the Conway normal form is
strictly positive, not that the surreal itself is positive. The element $p$ may
have either sign, and $p=0$ is included. In particular, $\J$ is an ordered
$\R$-vector space.

Our exponential predicates have the form
\begin{equation}\label{eq:introF}
 R_F(x_1,\ldots,x_d)\quad\Longleftrightarrow\quad
 \sum_{j=1}^{N}c_j\Exp\!\left(b_j+\sum_{k=1}^{d}\alpha_{jk}x_k\right)=0,
\end{equation}
where $c_j,b_j\in\Qb$ and $\alpha_{jk}\in\A$. Only the variables $x_k$ are
quantified, and they range over $\Oz$. The exponential values are not another
quantified sort. Products of quantified variables, nested exponentials, and order
comparisons between sums of exponentials are not part of the language.

The notation $\Exp$ in \eqref{eq:introF} requires no global surcomplex
exponential. Its domain is $\Qb+\J$, and the ordinary complex component is
algebraic. The precise construction appears in \cref{sec:foundations}.

\subsection{The main conclusions}
\paragraph{Exact algebraic dependence.}
For any $z_1,\ldots,z_d\in\Qb+\J$, all algebraic relations over $\Qb$ among
$\Exp(z_1),\ldots,\Exp(z_d)$ are generated by the integer additive relations
among the $z_k$. In particular,
\begin{equation}\label{eq:introtrdeg}
 \trdeg_{\Qb}\Qb\bigl(\Exp(z_1),\ldots,\Exp(z_d)\bigr)
 =\dim_{\Q}\spanQ\{z_1,\ldots,z_d\}.
\end{equation}
A stronger statement computes the intersection of this field with $\C$ and
establishes linear disjointness over that intersection
(\cref{thm:toric,thm:constants}).

\paragraph{Decidability with algebraic slopes.}
Let $\Lalg$ be ordered additive arithmetic on $\Oz$, with congruence predicates,
expanded by every predicate \eqref{eq:introF}. Its theory is decidable, uniformly
for effectively specified algebraic coefficients. After a definable coordinate
split, every formula becomes a finite union of products of a Presburger-definable
integer set and an algebraically semilinear set in $\J$
(\cref{thm:algdecision,thm:rectangles}).

\paragraph{All elementary substructures.}
A substructure $G\subseteq\Oz$ is elementary for $\Lalg$ if and only if
\begin{equation}\label{eq:introelementary}
 G=\Z+W\quad\text{for a nonzero $\A$-vector subspace }W\subseteq\J.
\end{equation}
Consequently, every positive purely infinite $u$ gives a countable elementary
substructure $\Z+\A u$. These are additive, relational structures, not exponential
subfields or omnific subrings (\cref{thm:elementary,cor:onescale}).

\paragraph{A sharp nonlinear obstruction.}
The same language becomes undecidable after adjoining the single ternary predicate
$\exp(xy)=\exp(z)$. This does not depend on an unresolved transcendence conjecture:
the new predicate recovers multiplication and the old language already defines
the ordinary integers (\cref{thm:undecidable}).

\paragraph{Exact computational content of one slope.}
For a language with just the additive structure and the relation
$\exp(y)=\exp(\alpha x)$, where $\alpha$ is a fixed ordinary real irrational,
its first-order theory has exactly the Turing degree of the rational cut of
$\alpha$. It is decidable exactly for computable slopes
(\cref{thm:degree}).

\subsection{Relation to existing work and the repository}\label{subsec:provenance}
The Conway normal form, Gonshor exponential, and classical
Lindemann--Weierstrass theorem are established inputs, not new claims of this
article \cite{gonshor,mantova,popescu}. Presburger quantifier elimination is also
classical \cite{cluckers}. The argument is a combination and application of these
inputs. In particular, the linear-independence lemma itself is a short transport
argument; the more substantial proposed contribution is the precise language
analysis, its definable splitting, and the classification of elementary cores.

The repository \texttt{VladimirReshetnikov/Surreal} was inspected at commit
\repoSHA. Its catalogue describes extensive work on omnific Diophantine geometry,
set-sized quotients, automorphisms, and initial subgroups \cite{repo-catalogue}.
The quotient report explicitly uses $\Oz=\Z\oplus\J$, under its notation for the
purely infinite ideal, and studies much stronger ring-theoretic structures
\cite{repo-quotients}. The exponential-automorphism report concerns field
endomorphisms compatible with exponentiation, not the additive relational reduct
studied here \cite{repo-exp}.

No advanced theorem from those unrefereed reports is needed in the proofs below.
This avoids making our results conditional on their pending review. The overlap
is foundational notation and the motivation to distinguish the ordinary constant
term from the purely infinite component. A targeted repository search for
``Presburger'' returned no items but explicitly marked its results incomplete;
that is \emph{not} evidence of absence. Targeted public searches also did not
identify the precise classification proved here. These bounded searches do not
establish priority, and the manuscript makes no exhaustive novelty claim.

\section{Normal forms, size conventions, and the exponential domain}\label{sec:foundations}
\subsection{Class-sized objects and finite assertions}
We use the standard class of surreal numbers, with each individual surreal a
set-sized object and each normal form having set-sized support. Statements about
$\Oz$ as a first-order structure are understood formula by formula. Thus
$G\preccurlyeq\Oz$ means that every fixed finite first-order formula, with a finite
tuple of parameters in $G$, has the same truth value in the two structures.
No set containing all surreals, no class-indexed sum, and no unrestricted
class-saturation principle is used.

All group algebras and field extensions in the algebraic arguments are formed from
set-sized subgroups, usually finitely generated ones. Transcendence degree is used
only for the resulting set-sized fields. The proofs involving the entire $\J$
are finite algebraic or first-order elimination arguments and therefore also give
schemes of ordinary set-theoretic assertions.

\subsection{The purely infinite component}
The normal form of a surreal is
\begin{equation}\label{eq:normal}
 x=\sum_{\gamma\in S}r_\gamma\omega^\gamma,
\end{equation}
where $S\subseteq\No$ is a set, reverse well ordered, and the nonzero coefficients
$r_\gamma$ are real. Uniqueness, termwise addition, convolution multiplication, and
comparison by the leading nonzero term are the facts needed below
\cite[Chapter~5]{gonshor}; see also \cite[Theorem~2.10]{mantova}.
Define
\begin{equation}
 \J=\{p\in\No:\supp(p)\subseteq\No_{>0}\},\qquad \Oz=\Z\oplus\J.
\end{equation}
The empty support represents zero. Finite unions of the relevant supports and the
usual normal-form multiplication show that $\J$ is an $\R$-vector space and is
closed under multiplication. Consequently, $\Oz$ is a subring of $\No$.
The coefficient of $\omega^0$ gives a ring homomorphism $\ct:\Oz\to\Z$.
Put $\pp(x)=x-\ct(x)$.

\begin{lemma}\label{lem:basic}
Every nonzero $p\in\J$ satisfies $|p|>r$ for every positive real $r$. In particular,
$\J\cap\C=\{0\}$ inside $\SC$, and the order on $\Oz$ is lexicographic:
\begin{equation}\label{eq:lex}
 n+p<m+q\quad\Longleftrightarrow\quad
 p<q\ \lor\ (p=q\ \land\ n<m).
\end{equation}
The least positive omnific integer is $1$.
\end{lemma}
\begin{proof}
The leading exponent of a nonzero purely infinite normal form is positive. Its
leading monomial dominates every positive real, and its sign is that of the
leading coefficient. Apply this to $p-q$ to obtain \eqref{eq:lex}. If a positive
omnific integer has nonzero purely infinite part, it is larger than every positive
integer; otherwise it is an ordinary positive integer. The assertion about $\C$
follows since a real surreal that is an ordinary complex number is real.
\end{proof}

\begin{proposition}[Division with remainder]\label{prop:division}
For every ordinary integer $m\geq1$ and $x\in\Oz$, there are unique $y\in\Oz$ and
$r\in\{0,1,\ldots,m-1\}$ such that $x=my+r$. Moreover,
\begin{equation}\label{eq:congruence}
 x\equiv y\pmod m\quad\Longleftrightarrow\quad
 \ct(x)\equiv\ct(y)\pmod m.
\end{equation}
\end{proposition}
\begin{proof}
For $x=n+p$, write $n=mq+r$ in $\Z$ and take $y=q+p/m$. The space $\J$ is divisible,
so this is an omnific integer. Applying $\ct$ to two candidate decompositions gives
ordinary uniqueness of $r$ and of the integer quotient; then the purely infinite
parts agree as well. The same calculation gives \eqref{eq:congruence}.
\end{proof}

Together with the ordered abelian group axioms, discreteness and division with
remainder are the standard axioms of a Presburger group. Thus $(\Oz,+,-,<,0,1,
(\equiv_m)_{m\geq2})$ is a model of the complete theory of the corresponding
structure on $\Z$ \cite{cluckers}. Congruence means divisibility in the additive
group. The map $\ct$ is not order preserving: $\omega-100$ is positive even though
its constant term is negative.

\subsection{The two transcendence inputs}
Let $\exp:\No\to\No_{>0}$ denote Gonshor's increasing exponential group
isomorphism, extending the real exponential. Write
$\M=\{\omega^\gamma:\gamma\in\No\}$ for the Conway monomial group.
We use the established identity
\begin{equation}\label{eq:Gonshor} \exp(\J)=\M. \end{equation}
This is Gonshor's theorem, not the assertion that $\exp(p)=\omega^p$.
See \cite[Theorems~10.8--10.13]{gonshor}, as restated in
\cite[Theorem~2.16]{mantova}. Different purely infinite arguments have different
monomial exponentials.

The classical Lindemann--Weierstrass theorem, in the form we use, says
\begin{equation}\label{eq:LW}
 \begin{gathered}
 a_1,\ldots,a_s\in\Qb\text{ distinct}\quad\Longrightarrow\\
 e^{a_1},\ldots,e^{a_s}\text{ are linearly independent over }\Qb.
 \end{gathered}
\end{equation}
The $a_j$ are ordinary complex algebraic numbers. A modern proof of this classical
result is given in \cite{popescu}. Neither \eqref{eq:Gonshor} nor \eqref{eq:LW} is
reproved here. All uses of them will be explicit.

\begin{lemma}[Complex coefficients on real monomials]\label{lem:monomial}
Distinct Conway monomials are linearly independent over $\C$ inside $\SC$.
\end{lemma}
\begin{proof}
Separate the real and imaginary parts of a finite purported dependence. Uniqueness
of the real Conway normal form makes both coefficients at every monomial zero.
\end{proof}

\subsection{A canonical restricted surcomplex exponential}
\begin{definition}\label{def:domain}
Set $\Ealg=\Qb+\J\subseteq\SC$. For its unique decomposition $z=a+p$, define
\begin{equation}\label{eq:Expdef} \Exp(z)=e^a\exp(p). \end{equation}
\end{definition}
Uniqueness follows from \cref{lem:basic}. This domain is an additive $\Q$-vector
space. It is not asserted to be a complex algebra or a ring: for example,
$i\cdot p$ need not belong to it for nonzero $p\in\J$. Its real part
$\A+\J$ is a subring, and on that real domain \eqref{eq:Expdef} is precisely the
restriction of Gonshor's exponential.

The map $\Exp$ is a homomorphism from the additive group to $\SC^\times$.
Only the ordinary complex exponential is used for imaginary parts, and those
imaginary parts are algebraic and finite. No sine or cosine at an infinite
imaginary argument is chosen. This limited surcomplex extension is sufficient for
all coefficients and shifts in \eqref{eq:introF}.

\section{Faithful exponential group algebras}\label{sec:faithful}
\begin{theorem}[Finite exponential independence]\label{thm:independence}
If $z_1,\ldots,z_s$ are distinct elements of $\Ealg$, then
$\Exp(z_1),\ldots,\Exp(z_s)$ are linearly independent over $\Qb$.
\end{theorem}
\begin{proof}
Write $z_j=a_j+p_j$ and suppose
$\sum_{j=1}^s c_j\Exp(z_j)=0$, where $c_j\in\Qb$. Group the finite sum according to
the distinct values of $p_j$. This gives
\begin{equation}\label{eq:grouped}
 \sum_{p}\left(\sum_{j:p_j=p}c_j e^{a_j}\right)\exp(p)=0.
\end{equation}
By \eqref{eq:Gonshor}, the outer factors $\exp(p)$ are distinct Conway monomials.
Their coefficients are ordinary complex numbers. Hence \cref{lem:monomial} gives
$\sum_{j:p_j=p}c_j e^{a_j}=0$ for every group. Within a fixed group the $a_j$ are
distinct, since the original $z_j$ were distinct. Apply \eqref{eq:LW} separately
to each group to obtain $c_j=0$ for every $j$.
\end{proof}

\begin{remark}
The outer sum in \eqref{eq:grouped} is finite. The $p_j$ may individually have
arbitrary set-sized reverse-well-ordered supports. No exchange of exponentiation
with an infinite sum is made in the proof; all such analytic-normal-form content
is confined to the imported theorem \eqref{eq:Gonshor}.
\end{remark}

\begin{corollary}\label{cor:injective}
The map $\Exp:\Ealg\to\SC^\times$ is injective. If $z\in\Ealg$ and
$\Exp(z)\in\Qb$, then $z=0$ and $\Exp(z)=1$.
\end{corollary}
\begin{proof}
Equality of two distinct exponential values, or a relation
$\Exp(z)-c\Exp(0)=0$ with $z\ne0$, contradicts \cref{thm:independence}.
\end{proof}

For a set-sized additive subgroup $H\subseteq\Ealg$, the group algebra $\Qb[H]$
consists of finite sums $\sum c_h[h]$, with $[h][k]=[h+k]$.
\begin{corollary}[Faithful group algebra]\label{cor:groupalgebra}
The map
\begin{equation}
 \Phi_H:\Qb[H]\longrightarrow\SC,\qquad
 \sum_h c_h[h]\longmapsto\sum_h c_h\Exp(h)
\end{equation}
is an injective $\Qb$-algebra homomorphism.
\end{corollary}
\begin{proof}
Multiplicativity follows from the additive exponential law. Injectivity is
\cref{thm:independence} after collecting equal group elements.
\end{proof}

\begin{example}
The pair $\omega,\sqrt2\,\omega$ is $\Q$-linearly independent even though its
members lie on one real-linear ray. Consequently their exponentials are
algebraically independent, even over $\C$, by the next section. The identity
$\exp(\sqrt2\,\omega)=\exp(\omega)^{\sqrt2}$ does not contradict this: a fixed
irrational real power is not a polynomial operation.
\end{example}

\section{Exact toric kernels and ordinary constants}\label{sec:toric}
\subsection{All finite algebraic relations}
Fix $z_1,\ldots,z_d\in\Ealg$ and put
\begin{equation}
 H=\sum_{k=1}^d\Z z_k,\qquad
 L=\{\ell\in\Z^d:\ell_1z_1+\cdots+\ell_dz_d=0\}.
\end{equation}
The group $H$ is finitely generated and torsion free. The relation lattice $L$ is
saturated: $m\ell\in L$ for a nonzero integer $m$ implies $\ell\in L$.

\begin{theorem}[Exact toric kernel]\label{thm:toric}
For the Laurent evaluation map
\begin{equation}
 \operatorname{ev}_z:\Qb[X_1^{\pm1},\ldots,X_d^{\pm1}]\to\SC,
 \qquad X_k\longmapsto\Exp(z_k),
\end{equation}
one has
\begin{equation}\label{eq:torickernel}
 \ker(\operatorname{ev}_z)=(X^\ell-1:\ell\in L).
\end{equation}
Its image is isomorphic to $\Qb[H]$. If
$r=\dim_\Q\spanQ\{z_1,\ldots,z_d\}$, then the image is a Laurent polynomial algebra
in $r$ variables, and
\begin{equation}\label{eq:trdeg}
 \trdeg_{\Qb}\Qb\bigl(\Exp(z_1),\ldots,\Exp(z_d)\bigr)=r.
\end{equation}
For ordinary polynomial evaluation, the kernel is generated by all
$X^u-X^v$ with $u,v\in\N^d$ and $u-v\in L$.
\end{theorem}
\begin{proof}
Send $X^u$ to $[\sum u_kz_k]$. This gives a surjection from the Laurent algebra to
$\Qb[H]$. Composing it with the injective map in \cref{cor:groupalgebra} is exactly
$\operatorname{ev}_z$. A Laurent polynomial $\sum c_uX^u$ lies in its kernel if and
only if the sum of its coefficients in every coset of $L$ is zero. In each finite
coset group, choose one exponent $u_0$. The corresponding polynomial is then a
linear combination of $X^u-X^{u_0}$, each of which is the Laurent unit $X^{u_0}$
times $X^{u-u_0}-1$. This proves \eqref{eq:torickernel}.

Choose a $\Z$-basis of the free group $H$. Its cardinality is $r$, and its group
algebra is $\Qb[T_1^{\pm1},\ldots,T_r^{\pm1}]$. Taking fraction fields gives
\eqref{eq:trdeg}. The same coefficient-grouping argument for a polynomial, using
only its nonnegative exponent vectors, proves the final assertion directly.
\end{proof}

\begin{corollary}[Independence criterion]\label{cor:algind}
A finite family in $\Ealg$ is $\Q$-linearly independent if and only if its
exponentials are algebraically independent over $\Qb$. In particular, this holds
for every finite family of omnific integers.
\end{corollary}

A $\Z$-basis of $L$ suffices to generate the Laurent ideal in
\eqref{eq:torickernel}: use $X^{u+v}-1=X^u(X^v-1)+(X^u-1)$ and the analogous
identity for $-u$. This makes the theorem constructive \emph{when a relation
lattice is supplied}. It is not an algorithm for discovering all rational
relations among arbitrary, uncoded surreal inputs.

\subsection{Exactly which standard complex numbers occur}
Let
\begin{equation}
 \begin{gathered}
 H_0=H\cap\C=H\cap\Qb,\\
 K=\Qb\bigl(\Exp(h):h\in H_0\bigr),\qquad
 F=\Qb\bigl(\Exp(h):h\in H\bigr).
 \end{gathered}
\end{equation}
The second description of $H_0$ follows from $\J\cap\C=0$.

\begin{theorem}[Constant field and linear disjointness]\label{thm:constants}
The fields $F$ and $\C$ are linearly disjoint over $K$, and
\begin{equation}\label{eq:constantfield} F\cap\C=K. \end{equation}
If $s$ is the $\Q$-rank of the purely infinite parts of $z_1,\ldots,z_d$, then
\begin{equation}
 \trdeg_{\C}\C\bigl(\Exp(z_1),\ldots,\Exp(z_d)\bigr)=s.
\end{equation}
\end{theorem}
\begin{proof}
The subgroup $H_0$ is saturated in $H$: if a nonzero integer multiple of $h$ has
zero purely infinite part, then $h$ has zero purely infinite part. Thus $H/H_0$ is
finitely generated torsion free, hence free, and the exact sequence splits.
Choose a basis of $H_0$ and elements $h_1,\ldots,h_s$ representing a basis of a
complement. The purely infinite parts $p_1,\ldots,p_s$ of these $h_j$ are
$\Q$-linearly independent. Put $t_j=\Exp(h_j)$.

For distinct integer exponent vectors $u$, the purely infinite elements
$\sum u_jp_j$ are distinct. A finite Laurent-polynomial relation in the $t_j$ with
complex coefficients therefore becomes a relation among distinct Conway
monomials, with each coefficient multiplied by a nonzero ordinary exponential.
By \cref{lem:monomial} all coefficients vanish. Hence $t_1,\ldots,t_s$ are
algebraically independent over $\C$, not merely over $K$.

We have $F=K(t_1,\ldots,t_s)$. The map
$\C\otimes_K K[t_1,\ldots,t_s]\to\C[t_1,\ldots,t_s]$ is an isomorphism, and
localizing at the nonzero polynomials over $K$ stays injective in
$\C(t_1,\ldots,t_s)$. This is linear disjointness. If $c\in F\cap\C$ but
$c\notin K$, the $K$-independent pair $1,c$ would be linearly dependent over $\C$,
a contradiction. The transcendence-degree statement follows from the same
algebraic independence and the chosen complement.
\end{proof}

\begin{corollary}[The omnific constant-field formula]\label{cor:ozconstants}
For $x_1,\ldots,x_d\in\Oz$, let $H=\sum\Z x_k$ and write
$H\cap\Z=g\Z$ with $g\geq0$. Then
\begin{equation}
 \Qb(\exp x_1,\ldots,\exp x_d)\cap\C
 =\begin{cases} \Qb,&g=0,\\ \Qb(e^g),&g>0. \end{cases}
\end{equation}
\end{corollary}

\begin{example}[A hidden ordinary exponential]
For $x_1=\omega$ and $x_2=2\omega+1$, the arguments are $\Q$-linearly independent,
so their exponentials are algebraically independent over $\Qb$. Nevertheless,
\begin{equation}
 \frac{\exp(2\omega+1)}{(\exp\omega)^2}=e.
\end{equation}
Their generated field is $\Qb(e,\exp\omega)$ and its intersection with $\C$ is
exactly $\Qb(e)$.
\end{example}

\begin{example}[An algebraic extension need not be present]
For $x_1=\omega$, $x_2=\omega+2$, the corresponding intersection is
$\Qb(e^2)$, not $\Qb(e)$. Indeed, $e$ is transcendental and
$T\notin\Qb(T^2)$: the automorphism $T\mapsto-T$ fixes the latter field but not
$T$. The existence of a square root in the ambient surcomplex field does not force
that root to belong to a specified subfield.
\end{example}

\begin{example}[Finite coefficients do not automatically become constants]
For $x_1=\omega+1$ and $x_2=\sqrt2\,\omega$, their purely infinite parts are
$\Q$-linearly independent. The ordinary complex constants in the generated field
are exactly $\Qb$, even though $\exp(x_1)=e\exp(\omega)$ visibly contains an ordinary
transcendental coefficient in its normal form.
\end{example}

\section{Balanced partitions eliminate finite exponential equations}\label{sec:partitions}
\begin{definition}
For coefficients $c=(c_1,\ldots,c_N)\in\Qb^N$, a partition $\mathcal P$ of
$\{1,\ldots,N\}$ is \emph{balanced} if $\sum_{j\in B}c_j=0$ for every block
$B\in\mathcal P$. Denote the finite set of balanced partitions by $\bal(c)$.
\end{definition}

\begin{theorem}[Balanced-partition elimination]\label{thm:partition}
For any $z_1,\ldots,z_N\in\Ealg$ and $c_1,\ldots,c_N\in\Qb$,
\begin{equation}\label{eq:partition}
 \sum_{j=1}^{N}c_j\Exp(z_j)=0
 \quad\Longleftrightarrow\quad
 \bigvee_{\mathcal P\in\bal(c)}
 \ \bigwedge_{B\in\mathcal P}\ \bigwedge_{j\in B}(z_j=z_{b(B)}),
\end{equation}
where $b(B)$ is a fixed representative of $B$. No inequalities between different
blocks are required.
\end{theorem}
\begin{proof}
Partition the indices according to their actual equal exponent values. By
\cref{thm:independence}, the exponential sum vanishes if and only if the
coefficient sum in every actual equality class is zero. This actual partition
then occurs on the right of \eqref{eq:partition}.

Conversely, if a right-hand conjunction holds, each block contributes zero:
\[\Exp(z_{b(B)})\sum_{j\in B}c_j=0.\]
Adding these finitely many zeros gives the
left side. Distinct chosen blocks may have equal exponent values; their zero
contributions simply coalesce. This explains why no cross-block disequality is
needed.
\end{proof}

For $N=0$ the sum is zero, and the right side is the true empty conjunction
associated with the empty partition. A nonempty list with no balanced partition
gives false. Zero coefficients may be deleted or allowed as singleton blocks.

\begin{corollary}[Affine exponential zero loci]\label{cor:affine}
The zero set in $\Oz^d$ of \eqref{eq:introF} is a finite union of sets defined by
finite systems of algebraic-affine linear equations. The union and those systems
are effectively obtainable from the algebraic coefficients, shifts, and slopes.
\end{corollary}
\begin{proof}
Every affine exponent lies in $\Ealg$. Apply \cref{thm:partition}, and replace each
within-block equality by
\begin{equation}\label{eq:lineareq}
 \sum_k(\alpha_{jk}-\alpha_{b(B)k})x_k+(b_j-b_{b(B)})=0.
\end{equation}
Algebraic addition and equality decide balance exactly. No exponential constant
has to be approximated numerically.
\end{proof}

\begin{remark}[An effective transcendence application]
Classical Lindemann--Weierstrass justifies a symbolic identity test. The algorithm
does not search for a numerical lower bound separating nearly equal exponential
sums. Even with ordinary complex exponentials, zero testing reduces to algebraic
coefficient sums and affine exponent equality.
\end{remark}

\subsection{Worked equations}
\begin{proposition}\label{prop:examples}
For omnific-integer variables the following equivalences hold:
\begin{align}
 \exp x+\exp y=2&\quad\Longleftrightarrow\quad x=y=0;\label{eq:sum2}\\
 \exp x+\exp y=\exp u+\exp v
 &\quad\Longleftrightarrow\quad
 (x=u\land y=v)\lor(x=v\land y=u);\label{eq:pairs}\\
 \exp(x+y)-\exp x-\exp y+1=0
 &\quad\Longleftrightarrow\quad x=0\lor y=0.\label{eq:factor}
\end{align}
More generally, equality of two finite sums of unit-coefficient exponentials is
equivalent to equality of the multisets of their exponents, including
multiplicities.
\end{proposition}
\begin{proof}
For \eqref{eq:sum2}, the coefficient list $(1,1,-2)$ has no proper nonempty
zero-sum subset; all three exponents $x,y,0$ must agree. For the multiset assertion,
group equal exponents and apply \cref{thm:independence}; the net integer
multiplicity in each group is zero. This gives \eqref{eq:pairs}. Finally, the left
side of \eqref{eq:factor} is $(\exp x-1)(\exp y-1)$ in a field. The factors vanish
exactly at zero by injectivity. Alternatively \cref{thm:partition} yields the
same alternatives.
\end{proof}

The conclusion in \eqref{eq:sum2} is special to the prescribed exponent domain;
it is false over all real arguments. A positive real $t\ne1$ with $t<2$ gives
real exponents $\log t,\log(2-t)$. Such logarithmic constants need not belong to
$\A+\J$.

\section{The languages and the rational-slope theorem}\label{sec:rational}
\subsection{Formal syntax and effective input}
The base language is
\begin{equation}
 \LPr=\{+,-,<,0,1,(\equiv_m)_{m\geq2}\}.
\end{equation}
All its function symbols act on the sole domain $\Oz$. The language $\Lalg$
adds one relation symbol $R_F$ for each finite specification
\eqref{eq:introF}; its coefficients and shifts are ordinary complex algebraic
numbers and its slopes are real algebraic numbers. The sublanguage $\Lrat$ uses
only rational slopes, while retaining algebraic coefficients and algebraic shifts.

An effective algebraic-number code consists, for example, of a defining rational
polynomial and a root-isolating interval or complex rectangle. Standard algebraic
number arithmetic decides validity, equality, field operations, realness, and the
sign of real algebraic numbers. A finite family can be embedded effectively in one
number field and expressed in a rational basis. These are the algebraic
computations used by the decision procedure. They can equivalently be included as
an exact algebraic-number subroutine specification; standard algorithms are treated in \cite{cohen}.

Ordinary integer parameters are terms in $1$. Arbitrary surreal constants are not
part of the effective input alphabet. Semantic results with arbitrary finite
parameter tuples will be stated separately; they do not imply an algorithm for
reading arbitrary surreal normal forms.

\subsection{Reduction to Presburger arithmetic}
\begin{lemma}\label{lem:rationalaffine}
For rational $q_1,\ldots,q_d$ and algebraic $b\in\Qb$, the relation
\begin{equation}
 q_1x_1+\cdots+q_dx_d+b=0,\qquad x_k\in\Oz,
\end{equation}
is either empty or is definable by one integer-affine equation in $\LPr$.
This distinction and the equation are effective.
\end{lemma}
\begin{proof}
The rational combination of the $x_k$ belongs to $\Q+\J$. If it equals the
ordinary algebraic number $-b$, its purely infinite part vanishes and $b$ is
rational. Thus for nonrational $b$ there are no solutions. If $b$ is rational,
clear all denominators. The resulting equation uses only integer multiples,
addition, and an ordinary integer constant, which are terms of $\LPr$.
\end{proof}

\begin{theorem}[Rational slopes add no Presburger expressive power]\label{thm:rational}
Every $\Lrat$-formula translates effectively to an equivalent $\LPr$-formula on
$\Oz$. Consequently its theory is decidable and
\begin{equation}\label{eq:Zelementary}
 (\Z,\Lrat)\preccurlyeq(\Oz,\Lrat),
\end{equation}
where the predicates on $\Z$ are obtained by the same exponential interpretation.
\end{theorem}
\begin{proof}
Use \cref{thm:partition} to eliminate every exponential predicate and
\cref{lem:rationalaffine} on each resulting affine equality. Leave the Boolean
connectives and quantifiers unchanged. This is a formula-by-formula translation,
so it applies to arbitrary alternations of quantifiers, not just existential
sentences.

The base structures on $\Z$ and $\Oz$ satisfy the complete Presburger theory.
The inclusion of $\Z$ preserves its quantifier-free language, including congruences
by \eqref{eq:congruence}. Presburger quantifier elimination therefore makes this
inclusion elementary. Applying the translation proves \eqref{eq:Zelementary}.
The same elimination decides sentences effectively \cite{cluckers}.
\end{proof}

\begin{remark}
The theory contains order \emph{on the omnific variables}, not order comparisons
between exponential sums. Nor does the use of a rational slope declare division
to be a total function on $\Oz$. Rational multiplication is used only in the
ambient affine exponent, and the final relations are expressed after clearing
denominators.
\end{remark}

\section{Algebraic slopes and effective two-sort elimination}\label{sec:algebraic}
\subsection{Splitting every linear equation}
\begin{lemma}[Algebraic coordinate splitting]\label{lem:split}
Let $\alpha_k\in\A$ and $b\in\Qb$. Write $x_k=n_k+p_k$ as in
\eqref{eq:introsplit}. Then
\begin{equation}\label{eq:splitlinear}
 \sum_k\alpha_kx_k+b=0
 \quad\Longleftrightarrow\quad
 \left(\sum_k\alpha_kp_k=0\right)
 \land\left(\sum_k\alpha_kn_k+b=0\right).
\end{equation}
The second conjunct is effectively equivalent to a finite system of integer-affine
equations in the $n_k$.
\end{lemma}
\begin{proof}
The first sum on the right is in $\J$; the second is an ordinary complex number.
Their intersection is zero, so the split is exact. Let $B$ be a number field
containing the finitely many $\alpha_k$ and $b$, and choose a rational basis
$\beta_1,\ldots,\beta_t$ of $B$. Write
$\alpha_k=\sum_h a_{hk}\beta_h$ and $b=\sum_h b_h\beta_h$ with rational
coordinates. The ordinary equation is equivalent to
$\sum_k a_{hk}n_k+b_h=0$ for every $h$. Clear denominators in each equation.
\end{proof}

Coefficients in a number-field basis are \emph{fixed rational numbers}; they are
not new quantified variables. This step introduces no multiplication of integer
variables.

\subsection{Quantifier elimination for the vector component}
Let $\Lvec$ be the language of ordered abelian groups with a unary scalar map
$\lambda_\alpha$ for every $\alpha\in\A$. The intended structures are nonzero
ordered $\A$-vector spaces. No scalar-field sort is quantified.

\begin{lemma}[Ordered vector-space elimination]\label{lem:vectorQE}
The theory of nonzero ordered $\A$-vector spaces has quantifier elimination and
is complete and decidable. Every order-preserving $\A$-linear embedding between
such spaces is elementary.
\end{lemma}
\begin{proof}
All terms are affine linear expressions in vector variables with fixed algebraic
scalar coefficients and, when parameters are allowed, fixed vector parameters.
To eliminate one existential variable, put the quantifier-free formula into a
finite disjunction of conjunctions. Replace a disequality by its two strict-order
alternatives, and a weak inequality by its equality and strict-order alternatives.
It is enough to eliminate $y$ from a conjunction of linear equalities and strict
inequalities.

If an equality has a nonzero coefficient of $y$, solve it uniquely for $y$ and
substitute in the other conditions. If no equality has such a coefficient,
retain the equalities independent of $y$. After dividing by nonzero coefficients
and reversing inequalities when their coefficients are negative, every remaining
condition involving $y$ is $a_i<y$ or $y<b_j$. Existence is equivalent to all
$a_i<b_j$ when both kinds of bound occur. For finitely many such bounds, take the
maximum lower bound and minimum upper bound; their midpoint is a witness. With
only lower bounds or only upper bounds, unboundedness gives a witness. With no
bounds, nonemptiness suffices. Ordered nonzero vector spaces are dense and have
no endpoints, since halving a positive difference and adding a positive vector
are available.

This proves effective elimination of one variable. Induction eliminates all
quantifiers. Scalar arithmetic, equality, and signs in $\A$ are decidable. A
parameter-free quantifier-free sentence has only the vector constant zero, so its
truth is independent of the particular nonzero ordered vector space. This gives
completeness and decidability. Finally, a linear order embedding preserves every
quantifier-free formula, hence every formula by elimination.
\end{proof}

This elementary proof is the linear case of the general ordered-vector-space
setting studied in the literature on linear o-minimal structures
\cite{loveys-peterzil}. It also explains why different vector-space dimensions or
Archimedean ranks need not be distinguished by first-order sentences in this
language: scalars are named operations, not quantified elements of a field sort.

\subsection{The translation theorem}
\begin{theorem}[Effective two-sort reduction]\label{thm:algdecision}
Every $\Lalg$-formula on $\Oz$ translates effectively, under $x_k=n_k+p_k$, into
a formula in the disjoint two-sorted structure
\begin{equation}\label{eq:twosort}
 (\Z,\LPr)\ \sqcup\ (\J,\Lvec).
\end{equation}
The theory of $(\Oz,\Lalg)$ is decidable. The translated formulas admit effective
quantifier elimination in the two-sorted language.
\end{theorem}
\begin{proof}
First replace each exponential atom by \cref{thm:partition}. Apply
\cref{lem:split} to the resulting algebraic-affine equalities. This yields
conjunctions of Presburger equations in $n$ and vector equations in $p$.

For base-language atoms, addition and negation split coordinatewise, constants
$0,1$ become $(0,0),(1,0)$, congruence is given by \eqref{eq:congruence}, and order
is given by \eqref{eq:lex}. For example,
\begin{equation}
 x_1+2x_2<3\quad\Longleftrightarrow\quad
 p_1+2p_2<0\ \lor\
 (p_1+2p_2=0\land n_1+2n_2<3).
\end{equation}
Equality splits into equality of both coordinates. Replace an existential
quantifier over $x$ by $\exists n\in\Z\,\exists p\in\J$, and a universal quantifier
by the corresponding pair of universal quantifiers. The bijection
\eqref{eq:introsplit} proves the translation correct by induction on formulas.

Every atomic formula of the result belongs to exactly one of the two sorts.
Finite Boolean combinations can be put into a finite disjunction of formulas
$A(n)\land B(p)$. To eliminate an integer variable, distribute its existential
quantifier over this disjunction and use Presburger elimination in $A$; $B$ does
not involve that variable. Eliminate a vector variable in the same way using
\cref{lem:vectorQE}. Universal quantifiers are handled by negation. This terminates
for any fixed finite input formula. For a sentence, the final truth value is
effectively determined. No bound on computational efficiency is asserted.
\end{proof}

\begin{example}[An irrational equation separates the coordinates]
The relation
\begin{equation}\label{eq:sqrt2example}
 \Exp(y)=\Exp(\sqrt2\,x),\qquad x,y\in\Oz,
\end{equation}
is equivalent to $y=\sqrt2\,x$. Its coordinate conditions are
$p_y=\sqrt2\,p_x$ and $n_y-\sqrt2\,n_x=0$. Expanding in the rational basis
$1,\sqrt2$ forces $n_y=n_x=0$. Thus \eqref{eq:sqrt2example} describes a scalar
map on purely infinite elements and cannot have a nonzero ordinary integer input.
This observation is the key to the next section.
\end{example}

\section{Definable coordinates and a normal form for all definable sets}\label{sec:definable}
\subsection{Recovering the integer and purely infinite sorts internally}
\begin{theorem}[Definable splitting]\label{thm:defsplit}
In $(\Oz,\Lalg)$ the subsets $\J$ and $\Z$, and the graph of $\ct$, are
parameter-free definable. Explicit definitions are
\begin{align}
 \Pure(x)&\quad:\Longleftrightarrow\quad
 \exists y\,\bigl(\Exp(y)=\Exp(\sqrt2\,x)\bigr),\label{eq:Pure}\\
 \Int(x)&\quad:\Longleftrightarrow\quad
 \forall p\,\bigl(\Pure(p)\land p>0\ \Longrightarrow\ -p<x<p\bigr),\label{eq:Int}\\
 \ctgraph(x,n)&\quad:\Longleftrightarrow\quad
 \Int(n)\land\Pure(x-n).\label{eq:CT}
\end{align}
For every $\alpha\in\A$, the scalar action on $\J$ is definable by
\begin{equation}\label{eq:scalardef}
 \Pure(p)\land\Pure(q)\land\bigl(\Exp(q)=\Exp(\alpha p)\bigr).
\end{equation}
Consequently the structure is definably equivalent to the two-sorted structure
\eqref{eq:twosort}, with reconstruction map $(n,p)\mapsto n+p$.
\end{theorem}
\begin{proof}
By injectivity, the witness in \eqref{eq:Pure} must be $y=\sqrt2\,x$. For
$x=n+p$, this is an omnific integer exactly when $\sqrt2\,n\in\Z$, which happens
exactly for $n=0$. Thus \eqref{eq:Pure} defines $\J$.

An ordinary integer is strictly between $-p$ and $p$ for every positive
$p\in\J$, by \cref{lem:basic}. Conversely, if $x=n+p_0$ and $p_0\ne0$, take
$p=|p_0|/2\in\J_{>0}$. The leading purely infinite comparison gives
$|x|>p$, so \eqref{eq:Int} fails. Hence it defines exactly $\Z$.
The unique decomposition \eqref{eq:introsplit} now proves \eqref{eq:CT}.
Equation \eqref{eq:scalardef} is the graph $q=\alpha p$, again by injectivity.

The induced ordered additive structure on the integer sort is Presburger
arithmetic, and on the pure sort we have all ordered $\A$-vector operations.
\Cref{thm:algdecision} defines every original relation from these two sorts.
Conversely \eqref{eq:Pure}--\eqref{eq:scalardef} and the inherited group operations
define the two sorts and their coordinate projections in the original structure.
The forward and reverse constructions compose via the definable bijection
$(n,p)\mapsto n+p$. This proves the stated equivalence, including the definable
identification of each composition with the original structure.
\end{proof}

\begin{remark}
Definability of ordinary $\Z$ does not by itself imply undecidability. Here its
induced structure contains addition and order but not multiplication. The explicit
reduction proves that no hidden integer multiplication is supplied by the
exponential-equality predicates.
\end{remark}

\subsection{Rectangular normal form and orthogonality}
Call a subset of $\J^d$ \emph{$\A$-semilinear} if it is a finite Boolean combination
of affine linear equalities and inequalities with scalar coefficients in $\A$
and fixed vector parameters in $\J$. Equivalently, it is a finite union of sets
specified by finite systems of such conditions.

\begin{theorem}[Rectangular normal form]\label{thm:rectangles}
Every $\Lalg$-definable set $X\subseteq\Oz^d$, allowing a finite tuple of parameters,
has a description
\begin{equation}\label{eq:rectangles}
 X=\bigcup_{j=1}^{r}
 \{n+p:n\in A_j\subseteq\Z^d,\ p\in B_j\subseteq\J^d\},
\end{equation}
where $A_j$ is Presburger definable and $B_j$ is $\A$-semilinear. The parameters
needed in the $A_j$ are ordinary integer coordinates of the original parameters;
the vector parameters needed in the $B_j$ are their purely infinite coordinates.
Conversely every set of the form \eqref{eq:rectangles} is $\Lalg$-definable.
\end{theorem}
\begin{proof}
Treat the original parameters as additional free variables during the translation
in \cref{thm:algdecision}. Quantifier elimination in the two sorts gives a finite
Boolean combination of sort-separated atoms. Put it into disjunctive normal form;
each conjunction is the intersection of one integer condition and one vector
condition. Presburger elimination and \cref{lem:vectorQE} give the indicated
forms. This proves \eqref{eq:rectangles} and the parameter statement. The converse
uses the definable coordinate maps and scalar graphs from \cref{thm:defsplit}.
\end{proof}

\begin{corollary}[No extra induced structure]\label{cor:stable}
The integer and purely infinite subsets are stably embedded: a definable subset
of either of their finite Cartesian powers can be defined using parameters from
that same subset. Their induced structures are exactly Presburger arithmetic and
ordered $\A$-vector-space geometry, respectively.
\end{corollary}
\begin{proof}
In \eqref{eq:rectangles}, a subset of $\Z^d$ is obtained by setting $p=0$. Each
$B_j$ either contains zero or does not, leaving a finite union of the corresponding
$A_j$. Similarly set $n=0$ for a subset of $\J^d$. The converse inclusions of the
induced languages were proved in \cref{thm:defsplit}.
\end{proof}

\begin{corollary}[Finite-image cross-component maps]\label{cor:finiteimage}
Every definable map from a definable subset of $\Z^m$ to $\J^n$ has finite image.
The same holds for a definable map from a subset of $\J^m$ to $\Z^n$, even with
parameters from $\Oz$.
\end{corollary}
\begin{proof}
The graph is a finite union of rectangles in the two sorts. A nonempty rectangle
$A\times B$ contained in a function graph can have at most one element in its
output factor $B$. Thus each nonempty rectangle contributes at most one output.
The proof in the reverse direction is identical.
\end{proof}

This is an exact separation statement, not merely a dimension comparison.
For example, there is no definable injection of the infinite set $\Z$ into $\J$
in this language, despite the existence of many set-theoretic injections and the
ambient operation $n\mapsto n\omega$. That operation would require multiplication
of a variable by the nonalgebraic parameter $\omega$, which is not one of the
allowed scalar maps on omnific variables.

\section{Elementary substructures, embeddings, and one-scale cores}\label{sec:cores}
\begin{theorem}[Classification of elementary substructures]\label{thm:elementary}
For a substructure $G\subseteq(\Oz,\Lalg)$ the following are equivalent:
\begin{enumerate}[label=(\roman*)]
\item $G\preccurlyeq(\Oz,\Lalg)$;
\item $G=\Z+W$ for a nonzero $\A$-vector subspace $W\subseteq\J$, equipped with its
inherited order and all inherited exponential-equality predicates.
\end{enumerate}
The assertion applies to set-sized $G$, and formula by formula to class-sized
substructures as well.
\end{theorem}
\begin{proof}
Suppose first that $G$ is elementary. It is an additive subgroup containing $1$,
so it contains every ordinary integer. If $x\in G$, its actual integer constant
term $n=\ct(x)$ is therefore in $G$, and $x-n\in G\cap\J$. Thus
$G=\Z+(G\cap\J)$. Set $W=G\cap\J$.

The sentence $\exists p\,(\Pure(p)\land p>0)$ holds in $\Oz$, so by elementarity
it holds in $G$. Since $\Pure$ is definable and elementarity preserves its truth
on parameters in $G$, this gives a nonzero member of $W$.
For $p\in W$ and fixed $\alpha\in\A$, the sentence with parameter $p$ asserting
existence of a pure $q$ satisfying \eqref{eq:scalardef} is true in $\Oz$. It is
therefore true in $G$. Its witness must equal the actual $\alpha p$ by injectivity
of $\Exp$. Hence $W$ is closed under every algebraic scalar, and is a nonzero
$\A$-vector subspace.

Conversely, let $G=\Z+W$ with $W\ne0$ an $\A$-vector subspace. It is a
substructure for the group operations and constants, with the relations inherited
from $\Oz$. The coordinate translation of \cref{thm:algdecision} works in $G$
with the two sorts $\Z$ and $W$: real algebraic combinations of the pure
coordinates remain in $W$, and lexicographic order and congruence have the same
formulas. By \cref{lem:vectorQE}, $W\preccurlyeq\J$ in the vector language. The
identity on the integer sort is elementary. Therefore every translated formula
with parameters from $\Z\sqcup W$ has the same truth value in
$\Z\sqcup W$ and $\Z\sqcup\J$. Translating back gives $G\preccurlyeq\Oz$.
\end{proof}

\begin{corollary}[One-scale elementary cores]\label{cor:onescale}
For every $u\in\J_{>0}$,
\begin{equation}\label{eq:onescale}
 G_u=\{n+a u:n\in\Z,\ a\in\A\}
 \preccurlyeq(\Oz,\Lalg).
\end{equation}
The set $G_u$ is countable. For positive $u,v\in\J$, the map
$n+a u\mapsto n+a v$ is an isomorphism of these structures.
\end{corollary}
\begin{proof}
The map $a\mapsto a u$ identifies the nonzero ordered $\A$-vector space $\A$ with
$\A u$. Apply \cref{thm:elementary}. Countability follows from that of $\Z$ and
$\A$. The displayed map is the identity on the integer coordinates and an ordered
$\A$-linear isomorphism on the pure coordinates, so the coordinate translation
preserves every relation and formula.
\end{proof}

For $u=\omega$, the set $G_\omega$ is not closed under multiplication:
$\omega^2\notin\Z+\A\omega$ by normal-form uniqueness. It is also not asserted to
contain its exponential values. Those values are used externally to interpret
relations, just as a relation on integers may be defined using a larger field.
This is why \eqref{eq:onescale} is compatible with the immense field-theoretic
complexity of the full surreal numbers.

\begin{corollary}[Classification of elementary embeddings]\label{cor:embeddings}
Let $G_W=\Z+W$ and $G_V=\Z+V$ be elementary substructures as above. Their elementary
embeddings are exactly the maps
\begin{equation}\label{eq:embeddingform}
 n+p\longmapsto n+T(p),
\end{equation}
where $T:W\to V$ is an order-preserving $\A$-linear embedding. In particular,
\begin{equation}
 \Aut(G_W,\Lalg)\cong\Aut_{\A,<}(W).
\end{equation}
\end{corollary}
\begin{proof}
An elementary embedding fixes ordinary integers, preserves the definable pure
subset and the graph of $\ct$, and preserves every scalar graph
\eqref{eq:scalardef}. Its restriction to $W$ is therefore an order-preserving
$\A$-linear embedding, and additivity gives \eqref{eq:embeddingform}. Conversely,
\cref{lem:vectorQE} makes every such $T$ elementary on vector coordinates.
Together with the identity on $\Z$, the two-sort translation proves elementarity
of \eqref{eq:embeddingform}. The automorphism statement is the bijective case.
\end{proof}

\begin{proposition}[The two slope languages genuinely differ]\label{prop:difference}
The inclusion $\Z\subseteq\Oz$ is elementary for $\Lrat$ but not for $\Lalg$.
\end{proposition}
\begin{proof}
The first assertion is \cref{thm:rational}. In $\Lalg$, the sentence
$\exists x\,(x>0\land\Pure(x))$ is true in $\Oz$. In the induced structure on
$\Z$, the defining equation $y=\sqrt2\,x$ has only the solution $x=y=0$, so the
sentence is false.
\end{proof}

\begin{remark}[Why one scale can be elementary]
A formula has only finitely many fixed algebraic scalar coefficients. It can
specify finitely many linear bounds and equations, but it cannot quantify over
all scalars to say that one positive vector is larger than every scalar multiple
of another. The vector-space elimination proof shows exactly why a nonempty
interval left by finitely many constraints already contains a witness in an
ordered algebraic-linear subspace. The ``one scale'' assertion is a
first-order-language phenomenon, not a claim that all surreal growth rates
collapse arithmetically.
\end{remark}


\section{One irrational slope and an exact computability classification}\label{sec:single}
The preceding arguments also identify the precise computational content of a
much smaller language. Here the slope may be an arbitrary ordinary real
irrational, not necessarily algebraic. This extension admits only its scalar
graph, not all exponential sums with that slope.

\begin{definition}
For a fixed $\alpha\in\R\setminus\Q$, let $\mathcal S_\alpha$ be the structure
on $\Oz$ in the finite signature $\{+,-,<,0,1,S\}$, where
\begin{equation}\label{eq:Salpha}
 S(x,y)\quad\Longleftrightarrow\quad
 \exp(y)=\exp(\alpha x).
\end{equation}
Put $K_\alpha=\Q(\alpha)$, as an ordered subfield of $\R$.
Congruence predicates can be added definitionally by
$x\equiv_m y\Longleftrightarrow\exists z\,(x-y=mz)$.
\end{definition}

The real argument $\alpha x$ is a legitimate surreal argument of Gonshor's
exponential even when its ordinary constant part is transcendental. Real
exponential injectivity gives $S(x,y)\Longleftrightarrow y=\alpha x$.
No use of Lindemann--Weierstrass is needed in this section.

\begin{theorem}[Single-slope decomposition]\label{thm:single}
The structure $\mathcal S_\alpha$ is definably equivalent to the disjoint union
of Presburger $\Z$ and the nonzero ordered $K_\alpha$-vector space $\J$.
Its elementary substructures are exactly
\begin{equation}\label{eq:singlecores}
 \Z+W,\qquad 0\ne W\subseteq\J\text{ a }K_\alpha\text{-vector subspace}.
\end{equation}
In particular, $\Z+K_\alpha u$ is a countable elementary core for every
$u\in\J_{>0}$.
\end{theorem}
\begin{proof}
Write $x=n+p$, $y=m+q$. Since $\alpha$ is irrational,
\begin{equation}\label{eq:Ssplit}
 S(x,y)\quad\Longleftrightarrow\quad n=m=0\ \land\ q=\alpha p.
\end{equation}
Indeed, the ordinary component equation $m=\alpha n$ forces $m=n=0$.
Thus $\exists y\,S(x,y)$ defines $\J$, and the formula \eqref{eq:Int}, with this
new pure-part definition, defines $\Z$. The decomposition and constant-term graph
are then definable as before.

On $\J$ the relation $S$ is the graph of scalar multiplication by $\alpha$.
Composition and addition give every polynomial in this operator with rational
coefficients. Rational scalar maps are definable by divisibility in the vector
group. For a nonzero polynomial value $g(\alpha)$, the graph of multiplication by
$f(\alpha)/g(\alpha)$ is defined by the equation
$g(\alpha)q=f(\alpha)p$ using these polynomial operator graphs. It has exactly one
solution $q$ for every $p$. Hence every scalar map of $K_\alpha$ is definable.

Conversely, the two coordinate sorts define $S$ by \eqref{eq:Ssplit} and the base
group language by lexicographic order and coordinatewise addition. The
ordered-vector-space elimination proof of \cref{lem:vectorQE} works verbatim over
any fixed ordered scalar field; effectiveness is not needed for this semantic
claim. The same arguments as in \cref{thm:elementary} therefore yield
\eqref{eq:singlecores}. The field $\Q(\alpha)$ is countable, which proves the last
assertion.
\end{proof}

For the computational statement, encode rational numbers and integer polynomials
by natural numbers. Let
\begin{equation}
 D_\alpha=\{q\in\Q:q<\alpha\},\qquad
 P_\alpha=\{f\in\Z[T]:f(\alpha)>0\}.
\end{equation}
For sets of finite codes, $A\leq_T B$ means that an algorithm with a membership
oracle for $B$ decides membership in $A$; $\equiv_T$ means reductions both ways.

\begin{theorem}[The theory has the degree of the slope]\label{thm:degree}
For every fixed ordinary real irrational $\alpha$,
\begin{equation}\label{eq:degree}
 \Th(\mathcal S_\alpha)\ \equiv_T\ P_\alpha\ \equiv_T\ D_\alpha.
\end{equation}
In particular, $\Th(\mathcal S_\alpha)$ is decidable if and only if $\alpha$ is a
computable real number. The reductions are statements about each fixed $\alpha$;
the final equivalence does not assert uniform zero testing from arbitrary names
of computable reals.
\end{theorem}
\begin{proof}
First, an oracle for $P_\alpha$ decides the signs and zero values of polynomials
at $\alpha$: zero is equivalent to neither $f$ nor $-f$ belonging to $P_\alpha$.
It therefore performs exact ordered-field arithmetic with rational-function
representatives in $K_\alpha$. Translate a sentence using \eqref{eq:Ssplit} and
eliminate its integer and vector quantifiers as before, consulting this oracle
for scalar equality, division, and sign. This proves
$\Th(\mathcal S_\alpha)\leq_T P_\alpha$.

In the reverse direction, write $f(T)=\sum_{j=0}^{d}a_jT^j$. Construct a sentence
asserting that there are $p_0,\ldots,p_d\in\J$ with $p_0>0$,
$S(p_j,p_{j+1})$ for $j<d$, and $\sum_j a_jp_j>0$. Here $\J$ is replaced by its
defining formula, and integer multiples are additive terms. The recurrence
forces $p_j=\alpha^jp_0$, so the sentence holds exactly when $f(\alpha)>0$.
This is an effective reduction $P_\alpha\leq_T\Th(\mathcal S_\alpha)$.

The reduction $D_\alpha\leq_T P_\alpha$ uses the linear polynomials $nT-m$ for
$q=m/n$ and $n>0$. To prove the converse for a fixed $\alpha$, distinguish two
cases. If $\alpha$ is algebraic, polynomial signs at it are decidable by exact
algebraic arithmetic, with an algebraic representation of this fixed number
built into the algorithm. If $\alpha$ is transcendental, every nonzero integer
polynomial has a nonzero value at $\alpha$. The cut oracle provides nested
rational intervals containing $\alpha$ with arbitrarily small widths. Evaluate
$f$ by rational interval arithmetic until its value interval excludes zero. This
terminates for every nonzero $f$; the zero polynomial is handled separately.
Thus $P_\alpha\leq_T D_\alpha$ in both cases.

Finally, for an irrational real $\alpha$, its rational cut is decidable precisely
when it is computable. A decidable cut gives rational approximations by bisection.
Conversely, computable approximations decide comparison with any fixed rational
by refinement, since equality with a rational is impossible. This completes the
proof.
\end{proof}

The countable elementary core in \cref{thm:single} also shows directly that the
first-order truth set in \eqref{eq:degree} is the theory of a set-sized structure.
There is no appeal to a global satisfaction class for a proper-class universe.

\begin{proposition}[Nonuniformity of arbitrary real-number names]\label{prop:nonuniform}
There is no single algorithm that, given a program computing rational
approximations with prescribed error bounds to an algebraic irrational $\alpha$,
and a sentence of the signature of $\mathcal S_\alpha$, always decides that
sentence. This remains false even when every supplied name is valid and every
slope is promised to be algebraic irrational.
\end{proposition}
\begin{proof}
For a machine index $e$, define
\begin{equation}
 \alpha_e=\begin{cases}
 \sqrt2+2^{-t},&\text{if machine }e\text{ first halts at step }t,\\
 \sqrt2,&\text{if it never halts}.
 \end{cases}
\end{equation}
There is a uniformly computable approximation name for $\alpha_e$: simulate the
machine sufficiently many steps for the requested precision; if it has not
halted, the effect of any later halting is at most the corresponding power of
$2^{-1}$. Combine this bound with a rational approximation to $\sqrt2$.
Every $\alpha_e$ is algebraic irrational, irrespective of whether the machine
halts.

A fixed sentence of the scalar-graph language expresses
$\alpha_e^2>2$: choose pure $p>0$, apply $S$ twice to obtain $q=\alpha_e^2p$, and
require $q>2p$. The sentence is true exactly when the machine halts. A uniform
decision procedure from those approximation names would therefore decide the
halting problem, which is impossible. The argument does not supply a minimal
polynomial or an isolating-root representation of $\alpha_e$ uniformly; those
are the stronger algebraic codes required by our earlier positive algorithm.
\end{proof}

This gives a second boundary, independent of nonlinear exponentiation. Even a
single partial scalar graph can carry noncomputable information when its real
slope is not effectively known. For algebraic slopes with the explicit algebraic
codes specified in \cref{sec:rational}, this obstruction is absent.

\section{Boundaries: tails, constants, and a single undecidable extension}\label{sec:boundaries}
\subsection{Allowing arbitrary finite constants}
The algebraicity assumption on finite exponent constants is essential to
\cref{thm:independence}. Allowing $\log2$ immediately gives
\begin{equation} \exp(\log2)-2\exp(0)=0. \end{equation}
The arguments are distinct. This does not contradict Lindemann--Weierstrass:
$\log2$ is not algebraic, since a nonzero algebraic argument cannot have the
algebraic exponential $2$. The failure is located at an explicitly excluded
finite constant.

\subsection{Allowing infinitesimal tails}
\begin{proposition}[A tail obstruction]\label{prop:tails}
Finite exponential independence over $\Qb$ fails on the full surreal field even
when the coefficients are integers.
\end{proposition}
\begin{proof}
Take $\varepsilon=\omega^{-1}$, $a=\log(1+\varepsilon)$, and
$b=\log\varepsilon$. Then $a>0$, $b<0$, and both differ from zero. The
exponential inverse law in $\No$ gives
\begin{equation}\label{eq:tailcounterexample}
 \exp(a)-\exp(b)-\exp(0)=(1+\varepsilon)-\varepsilon-1=0.
\end{equation}
The element $a$ is a nonzero infinitesimal, so it does not lie in $\A+\J$.
Indeed, for every positive real $r$, one has $1+\varepsilon<e^r$, so
$0<a<r$ by monotonicity. By \eqref{eq:Gonshor}, $b$, as the logarithm of a
monomial, is purely infinite. Thus one infinitesimal-tail argument already
produces a nontrivial three-term relation.
\end{proof}

The proposition does not claim that $\Ealg$ is maximal among all additive domains
with finite exponential independence. Maximality is a separate extension problem;
see \cref{q:extension}.

\subsection{Nonlinear exponents recover multiplication}
\begin{theorem}[A single nonlinear relation destroys decidability]\label{thm:undecidable}
Expand $(\Oz,\Lalg)$ by the ternary predicate
\begin{equation}\label{eq:Mul}
 \mathsf M(x,y,z)\quad\Longleftrightarrow\quad\exp(xy)=\exp(z).
\end{equation}
The first-order theory of the expansion is undecidable.
\end{theorem}
\begin{proof}
The product $xy$ is an omnific integer because $\Oz$ is a ring. Injectivity of the
real exponential makes \eqref{eq:Mul} equivalent to $xy=z$. Thus the new relation
is the graph of multiplication on $\Oz$.

The old formula \eqref{eq:Int} defines exactly the ordinary integers. Relativize
all quantifiers in an ordinary arithmetic sentence to $\Int$, replace each
multiplication term by its graph using $\mathsf M$, and retain ordinary addition,
order, $0$, and $1$. This is an effective interpretation of the true first-order
arithmetic of $\Z$, equivalently of $\N$ after restricting to nonnegative
integers. That theory is undecidable by the classical incompleteness and
undecidability results for arithmetic \cite{salehi,kurahashi}. A decision procedure
for the expansion would decide it, a contradiction.
\end{proof}

This is a statement about the \emph{full first-order theory}. The relativization
uses a quantified definition of $\Z$; no claim about an existential-only reduction
is being silently made. Likewise, this argument is not a claim that every
nonlinear exponential predicate leads to undecidability.

\subsection{What the decision theorem does not cover}
The language analyzed here does not contain multiplication of two variable
omnific integers, the graph of $x\mapsto\exp x$ as an internal function, arbitrary
surreal scalar multiplication, arbitrary real or complex coefficient constants,
or exponential inequalities. It also has no complex-variable quantifiers and no
exponential at an infinitely large imaginary argument.

In particular, allowing an inequality $\alpha x<y$ with an irrational algebraic
$\alpha$ on the \emph{finite integer coordinates} is a different extension: it
creates integer cuts involving irrational slopes rather than the rational linear
equalities produced in \cref{lem:split}. The current proof does not establish a
decision procedure for that extension. It would be incorrect to infer one merely
because algebraic slopes are allowed inside \emph{equality} predicates.

\section{Beyond algebraic coefficients: exact fibers for real-base equations}\label{sec:fibers}
There is a useful extension of the grouping argument that does not require
Lindemann--Weierstrass. It gives geometry over the purely infinite coordinates but
leaves a genuinely ordinary exponential equation on the integer coordinates.

\begin{theorem}[Ordinary-fiber criterion]\label{thm:fibers}
Let $c_j\in\C$, $\lambda_{jk}\in\R$, and consider
\begin{equation}\label{eq:realF}
 F(x)=\sum_{j=1}^{N}c_j\exp\!\left(\sum_k\lambda_{jk}x_k\right),
 \qquad x\in\Oz^d.
\end{equation}
For a fixed $n\in\Z^d$, put
\begin{equation}
 d_j(n)=c_j e^{\sum_k\lambda_{jk}n_k},\qquad
 P_j(p)=\sum_k\lambda_{jk}p_k.
\end{equation}
Then $F(n+p)=0$ if and only if there exists a partition $\mathcal P$ of
$\{1,\ldots,N\}$ such that for every block $B$,
\begin{equation}\label{eq:fiberblock}
 \sum_{j\in B}d_j(n)=0,
 \qquad P_j(p)=P_{b(B)}(p)\quad(j\in B).
\end{equation}
Thus every fixed integer fiber is a finite union of real-linear subspaces of
$\J^d$. No uniform decision procedure for which partitions satisfy the ordinary
coefficient equations is asserted.
\end{theorem}
\begin{proof}
Rewrite \eqref{eq:realF} as $\sum_jd_j(n)\exp(P_j(p))$. Since each $P_j(p)\in\J$,
distinct values produce distinct Conway monomials by \eqref{eq:Gonshor}.
Their complex coefficients vanish separately by \cref{lem:monomial}. The
balanced-partition argument now applies to these fixed complex numbers $d_j(n)$
without any use of their algebraicity. Within-block equality is real-linear and
homogeneous in $p$, which gives the fiber description.
\end{proof}

\begin{corollary}[Positive real-base ternary equations]\label{cor:ternary}
Fix ordinary real numbers $a,b,c>1$. All solutions in $\Oz^3$ to
\begin{equation}\label{eq:abc}
 a^x+b^y=c^z,\qquad a^x:=\exp((\log a)x),
\end{equation}
are exactly
\begin{equation}\label{eq:abcparam}
 x=n+\frac{\lambda}{\log a},\qquad
 y=m+\frac{\lambda}{\log b},\qquad
 z=k+\frac{\lambda}{\log c},
\end{equation}
where $n,m,k\in\Z$ satisfy $a^n+b^m=c^k$ and $\lambda\in\J$ is arbitrary.
\end{corollary}
\begin{proof}
In a fixed integer fiber the three nonzero coefficients have signs $+,+,-$.
A partition into at least two blocks would have a block consisting only of positive
coefficients or the lone negative coefficient, so it cannot be balanced. Hence all
three purely infinite exponents agree:
$(\log a)\pp(x)=(\log b)\pp(y)=(\log c)\pp(z)=\lambda$.
Their common coefficient equation is the ordinary equation
$a^n+b^m-c^k=0$. This proves necessity and gives \eqref{eq:abcparam}. Substitution
proves sufficiency. Division by the fixed nonzero real logarithms preserves $\J$.
\end{proof}

\begin{example}
Since $2^1+3^1=5^1$, every $\lambda\in\J$ gives the solution
\begin{equation}
 2^{\,1+\lambda/\log2}+3^{\,1+\lambda/\log3}
 =5^{\,1+\lambda/\log5}.
\end{equation}
Each exponent is an omnific integer. For $\lambda>0$ they are all positive
infinite. Their ordinary constant terms are nevertheless exactly $1$.
\end{example}

\begin{corollary}[An inhomogeneous anchor removes all infinite solutions]
For fixed ordinary $a,b>1$, every omnific-integer solution of
\begin{equation}\label{eq:anchor} a^x-b^y=1 \end{equation}
has $x,y\in\Z$. Thus \eqref{eq:anchor} has exactly its ordinary integer solutions.
\end{corollary}
\begin{proof}
The three coefficients in each integer fiber have signs $+,-,-$, so again only
the one-block partition can be balanced. The third purely infinite exponent is
zero because it represents the constant $1$. Therefore
$(\log a)\pp(x)=(\log b)\pp(y)=0$, whence both purely infinite parts vanish.
\end{proof}

These are transfer and parameterization theorems. They do not solve the ordinary
integer equations left on the base. In the homogeneous equation
\eqref{eq:abc}, positivity of an infinite omnific variable does not force its
ordinary constant term to be positive, so restrictions on $n,m,k$ must not be
inferred from order conditions on $x,y,z$ without using \eqref{eq:lex}.

\section{Effective procedures and finite verification}\label{sec:verification}
\subsection{A complete mathematical decision pipeline}
The proofs provide the following terminating algorithm for a finite $\Lalg$
sentence. This is a mathematical specification of the full procedure, not a claim
that the accompanying small program implements all its stages.
\begin{enumerate}
\item Normalize every relation atom to a finite exponential sum. Products of
exponential factors with allowed affine exponents can be flattened using
$\Exp(u)\Exp(v)=\Exp(u+v)$; Laurent powers also cause no difficulty. Remove zero
coefficients when convenient.
\item Enumerate the set partitions of the finite list of terms. Keep those with
zero algebraic coefficient sum in every block. For each retained partition,
output within-block affine exponent equalities.
\item Introduce integer and purely infinite coordinates for every omnific
variable. Use \cref{lem:split} and rational-basis arithmetic in the finite number
fields occurring in the input. Translate order and congruence by
\eqref{eq:lex} and \eqref{eq:congruence}.
\item Eliminate quantifiers in the two disjoint sorts, using Presburger
elimination and the ordered-vector-space procedure in \cref{lem:vectorQE}.
Evaluate the resulting closed Boolean expression.
\end{enumerate}
This deliberately simple partition algorithm may produce $B_N$ branches, where
$B_N$ is the $N$th Bell number, before accounting for later Boolean and Presburger
expansion. No polynomial-time or elementary-time complexity bound is claimed.

The toric-kernel calculation is a different interface: it starts with a finite
argument family together with an exact description of its additive relation
lattice. Given that lattice, basis reduction supplies finitely many Laurent
binomial generators. Arbitrary surreal input equality is not solved by treating a
finite coefficient specification as if it were a full normal-form oracle.

\subsection{What the supplied program checks}
The file \texttt{code/checks.py} uses exact rational arithmetic and exact arithmetic
in $\Q(\sqrt2)$. It contains a reusable finite balanced-partition generator and
checks the following layers independently of numerical exponential evaluation:
\begin{enumerate}
\item balanced partitions versus direct coefficient grouping for exhaustive small
coefficient and exponent-label lists, including empty and zero-coefficient cases;
\item the finite-support forms of \eqref{eq:sum2}, \eqref{eq:pairs}, and
\eqref{eq:factor};
\item the $\sqrt2$ pure-part detector and the coefficientwise splitting of
algebraic linear equations, using exact quadratic-field coordinates;
\item the inhomogeneous and homogeneous fiber criterion with positive and negative
ordinary coefficients at representative formal scales.
\end{enumerate}
The recorded run passed 22,074 finite cases, including 7,381 exhaustive
partition cases, 6,561 two-term multiset cases, and 2,000 quadratic-coordinate
splitting cases. The run is deterministic, uses no network access, and writes
its actual counts and outcomes to \texttt{check\_results.json}.
Exponential values in these checks are
represented in a finite formal group algebra. Therefore the tests check the
combinatorial and coordinate implementations, not Gonshor's theorem or
Lindemann--Weierstrass. They do not formally verify all quantifiers or all
proper-class assertions of the article.

\subsection{A formalization path}
A useful formalization starts with an abstract additive group $E$, a field $K$,
and an injective group-algebra evaluation into a field extension. The toric and
partition theorems then have no surreal prerequisites. A second layer formalizes
an ordered algebraic-vector space $V$ and the lexicographic group $\Z\oplus V$,
with finite algebraic-linear relation predicates. The splitting, rectangular
normal form, and elementary-substructure arguments belong to this layer.
Only the final bridge needs the concrete surreal normal form and Gonshor
exponential. This separation mirrors the distinction in the repository between
generic checked algebra and its concrete surreal applications
\cite{repo-catalogue,repo-exp}.

\section{Further research questions}\label{sec:questions}
The questions below are proposed continuations of this manuscript. They are not
presented as a catalogue of previously named open problems, and no priority claim
is attached to their formulation.

\begin{question}[Exponential inequalities]\label{q:inequality}
Which expansions by signs or order comparisons of finite exponential sums remain
decidable? A first target is rational slopes with real algebraic coefficients.
One must account for the leading monomial and the sign of its ordinary
exponential coefficient, not merely its vanishing. With algebraic irrational
slopes, clarify exactly which induced integer cuts appear.
\end{question}

\begin{question}[Countably many scalar operators]
\Cref{sec:single} classifies the structure with one irrational scalar graph.
For an effectively indexed countable family of named real slopes containing an
irrational slope, determine the oracle degree of the resulting theory in terms
of the ordered field they generate. Distinguish uniform approximation of the
slopes from uniform knowledge of all finite polynomial relations among them.
For finite families, also specify which extra presentation data make the
reductions uniform. The nonuniformity in \cref{prop:nonuniform} shows why a list
of approximation programs is not automatically enough.
\end{question}

\begin{question}[Extensions of the independent exponential domain]\label{q:extension}
Classify additive domains $E\supseteq\Qb+\J$ on which an extension of the
exponential group-algebra map over $\Qb$ remains injective. The counterexample
\eqref{eq:tailcounterexample} shows that arbitrary infinitesimal tails cannot all
be admitted. It does not exclude carefully chosen transcendental constants or
special tail subspaces. A useful objective is a criterion in terms of the
monomial and unit components of their exponentials; the extension of the
exponential itself must also be specified.
\end{question}

\begin{question}[More general ordinary coefficient fields]
For an effectively presented field of ordinary constants larger than $\Qb$,
identify sufficient independence and zero-testing hypotheses for the same
partition reduction. Distinguish unconditional transfers from those requiring
new ordinary transcendence information. Arbitrary real-base equations in
\cref{sec:fibers} provide a test case where only the purely infinite geometry is
automatic.
\end{question}

\begin{question}[The nonlinear boundary below multiplication]
The single ternary predicate in \cref{thm:undecidable} recovers multiplication.
Determine which smaller nonlinear families avoid this. Unary nonlinear predicates,
restricted compositions, or predicates whose nonlinearity is confined to a known
pure subspace are concrete candidates. Every proposal must be tested against
definable recovery of integer multiplication.
\end{question}

\begin{question}[Parameter-uniform effective theories]
Develop a computable presentation of selected parameter families in $\J$,
including an algorithm for the finite algebraic-linear order relations that the
vector-space elimination procedure asks. Determine when the elementary diagram of
such a parameter expansion is decidable. The arbitrary-parameter semantic theorem
in \cref{thm:rectangles} alone is not such an algorithm.
\end{question}

\begin{question}[Complexity and compact certificates]
Replace exhaustive balanced-partition enumeration by zero-sum circuit methods,
matroid data, or shared Boolean representations. Give upper and lower complexity
bounds for fixed arity, fixed coefficient field, and bounded numbers of
exponential terms. A certificate should record both the coefficient balance and
the affine equalities, so that it is independently checkable.
\end{question}

\begin{question}[Toric and arithmetic families]
Study uniform families of the constant-field theorem as the additive relation
lattice changes. Which ordinary constant fields can occur over prescribed pure
parts? How do finite-index changes such as $\Qb(e^2)$ versus $\Qb(e)$ interact with
specialization and normalization? The exact intersection formula is a starting
point, not a classification of all larger exponential subfields.
\end{question}

\begin{question}[Model-theoretic refinements]
Use the definable product description to study types, definable groups, canonical
parameters, and elementary embeddings in countable models of the resulting
complete theory. Relate the classification in \cref{thm:elementary} to the
repository's much stronger ring- and field-embedding questions without confusing
their languages. In particular, characterize additional predicates that first
make a one-scale core non-elementary.
\end{question}

\begin{question}[A larger canonical surcomplex domain]
Find useful, explicitly specified domains permitting more imaginary components
while retaining a canonical exponential and a controlled kernel. The present
domain avoids all infinite imaginary phases. A broader construction must state
its phase convention and track period relations before an exponential
independence or decidability theorem can be formulated correctly.
\end{question}

\begin{question}[Machine-checked bridge]
Formalize the abstract group-algebra and coordinate-splitting theorems first, then
supply the concrete bridge from Conway normal forms to
$\exp(\J)=\M$. Record exactly which parts are proved about an abstract model and
which are instantiated for the actual surreal field. This is a realistic route to
checking the entire dependency chain rather than only finite examples.
\end{question}

\section{Conclusion}
The exponential relations studied here have an exact algebraic explanation:
monomial separation handles the purely infinite component, and
Lindemann--Weierstrass handles the algebraic ordinary component. Every finite
relation then reduces to balanced additive coincidences of its exponents.

The resulting first-order theory is stronger than Presburger arithmetic when
algebraic irrational slopes are named, but it still has a precise decidable
structure: ordinary additive integers and a nonzero ordered algebraic-vector
space, with no additional interaction after definable splitting. This yields the
rectangular description of definable sets and the full classification of
its elementary substructures. The countable one-scale cores are consequences of
that classification, not approximations to the full surreal field. In the
single-irrational-slope language, the same splitting identifies the exact Turing
degree of the theory with the degree of the slope, while distinguishing
fixed-slope decidability from uniform decision using arbitrary approximation
names.

The nonlinear extension and the explicit tail counterexample locate real limits
of the results. The appropriate conclusion is therefore not a resolution of
unrestricted surreal exponentiation, but a proof-bearing structural and
algorithmic theory of a substantial, precisely bounded exponential-equality
fragment. Independent mathematical review and a broader priority search remain
necessary before treating the proposed contributions as established new results.

\appendix
\section{Proof dependencies and audit checklist}\label{app:audit}
\subsection{Dependencies}
The dependency chain is deliberately short. Normal-form uniqueness and
$\exp(\J)=\M$, together with ordinary Lindemann--Weierstrass, imply
\cref{thm:independence}. That theorem implies the exact toric kernel and the
balanced-partition criterion. The constant-field theorem needs only monomial
independence over $\C$ once the pure-part complement is chosen. Balanced partitions
and algebraic coordinate splitting reduce the logical language to Presburger
arithmetic and ordered-vector-space linear geometry. The latter reduction yields
decidability, definability of the two components, rectangular normal form, and the
elementary-core classification. The single-slope decomposition uses only
exponential injectivity and the same coordinate geometry, not
Lindemann--Weierstrass. Its degree calculation is proved by explicit oracle
reductions. Classical undecidability of integer arithmetic and of halting is
used only for the two negative boundary results.

No conjectural assertion in the source repository is an input. Nor are
Berarducci--Mantova differentiation, saturation, strong summability, set-sized
quotient rigidity, or a global complex exponential needed.

\subsection{Checks against likely overstatements}
\begin{longtable}{>{\raggedright\arraybackslash}p{0.31\textwidth}>{\raggedright\arraybackslash}p{0.63\textwidth}}
\toprule
Potential misreading & Actual assertion and safeguard\\
\midrule
\endhead
$\exp(p)=\omega^p$ & Not used. Only the set equality $\exp(\J)=\M$ and injectivity
are imported from Gonshor.\\[4pt]
A new classical transcendence theorem & Not claimed. Classical
Lindemann--Weierstrass is an explicit input; the finite independence argument is
its monomial transport.\\[4pt]
A global surcomplex exponential & Not defined. $\Exp$ acts only on $\Qb+\J$,
with ordinary finite algebraic imaginary part.\\[4pt]
A decision procedure for arbitrary surreal input & Not claimed. The effective
input is a finite formula with algebraic coefficient codes. Arbitrary surreal
parameters require additional effective data.\\[4pt]
Full exponential-field decidability & Not claimed. Only externally interpreted
finite-sum zero predicates with affine exponents are admitted.\\[4pt]
All algebraic-slope inequalities & Not admitted. Base order compares additive
omnific terms; the algebraic slopes occur inside exponential equalities.\\[4pt]
Quantifier elimination in the original one-sort signature & Not asserted. The
proved elimination is after a definable two-sort coordinate expansion.\\[4pt]
Uniform decisions from arbitrary real approximation names & Not claimed.
The fixed-slope Turing reductions need not be uniform in such names;
\cref{prop:nonuniform} proves the obstruction.\\[4pt]
A one-scale surreal subfield & Not claimed. $\Z+\A u$ is an additive relational
substructure; it is generally not a subring.\\[4pt]
Existential undecidability of the nonlinear extension & Not claimed. The proof
interprets full first-order arithmetic using a quantified integer definition.\\[4pt]
Verified novelty or formal proof & Neither is claimed. The literature search is
bounded, the repository search was incomplete, and finite tests are not a Lean
formalization.\\
\bottomrule
\end{longtable}

\section{Reproducibility and source record}\label{app:sources}
The accompanying archive contains this source, its compiled PDF,
\texttt{code/checks.py}, its recorded JSON output, a source-and-claim audit, and
build instructions. Build with \texttt{latexmk -pdf article.tex}; the bibliography
is contained in this file, so no external \texttt{.bib} file is needed. The checks
use Python's standard library only.

The public source inspection was performed on 23 September 2026. In the
Mantova--Matusinski PDF, the normal-form statement and the exponential monomial
statement were visually checked on PDF pages 7 and 11, respectively; the latter is
Theorem~2.16. In Cluckers' preprint, the Presburger language and the recalled
quantifier-elimination and decidability statements were visually checked on PDF
page 2. Popescu's version~2 abstract states exactly the complex
Lindemann--Weierstrass input used here. These checks establish what the references
say, not the novelty of the deductions made from them.

The repository catalogue, the exponential-automorphism report overview, and
relevant introductory portions of the quotient report were inspected through
the GitHub connector at the pinned commit. Its full collection is large; this was a focused
comparison of relevant topics, not a claim-by-claim audit of every repository
manuscript. The bibliography points to immutable commit URLs for the inspected
repository sources.

\begin{thebibliography}{99}
\bibitem{gonshor}
H. Gonshor, \emph{An Introduction to the Theory of Surreal Numbers},
London Mathematical Society Lecture Note Series 123, Cambridge University Press,
1986. Chapters 5, 8, and 10. DOI:
\href{https://doi.org/10.1017/CBO9780511629143}{10.1017/CBO9780511629143}.

\bibitem{mantova}
V. Mantova and M. Matusinski, Surreal numbers with derivation, Hardy fields and
transseries: a survey, \emph{Contemporary Mathematics} \textbf{697} (2017),
265--290. DOI: \href{https://doi.org/10.1090/conm/697/14057}{10.1090/conm/697/14057}.
Preprint \href{https://arxiv.org/abs/1608.03413}{arXiv:1608.03413};
Theorems 2.10 and 2.16 in the inspected PDF.

\bibitem{popescu}
S. A. Popescu, \emph{A simple and self-contained proof for the
Lindemann--Weierstrass theorem}, 2023, version 2, 17 September 2023.
\href{https://arxiv.org/abs/2306.14352v2}{arXiv:2306.14352v2}.
Used for an exposition of the classical theorem, not as a new attribution of it.

\bibitem{cluckers}
R. Cluckers, Presburger sets and $p$-minimal fields,
\emph{Journal of Symbolic Logic} \textbf{68}(1) (2003), 153--162.
DOI: \href{https://doi.org/10.2178/jsl/1045861509}{10.2178/jsl/1045861509}.
Preprint \href{https://arxiv.org/abs/math/0206197}{arXiv:math/0206197}, Section 1.1.

\bibitem{loveys-peterzil}
J. Loveys and Y. Peterzil, Linear o-minimal structures,
\emph{Israel Journal of Mathematics} \textbf{81}(1--2) (1993), 1--30.
DOI: \href{https://doi.org/10.1007/BF02761295}{10.1007/BF02761295}.
The particular ordered-vector-space elimination used here is proved in
\cref{lem:vectorQE}.

\bibitem{cohen}
H. Cohen, \emph{A Course in Computational Algebraic Number Theory},
Graduate Texts in Mathematics 138, Springer, 1993.
DOI: \href{https://doi.org/10.1007/978-3-662-02945-9}{10.1007/978-3-662-02945-9}.

\bibitem{salehi}
S. Salehi, \emph{Computation in Logic and Logic in Computation}, 2016.
\href{https://arxiv.org/abs/1612.06526}{arXiv:1612.06526}.
Includes the standard distinction between decidable additive theories and
undecidable integer arithmetic with addition and multiplication.

\bibitem{kurahashi}
T. Kurahashi, \emph{Incompleteness and undecidability of theories consistent with
$\mathsf R$}, 2022.
\href{https://arxiv.org/abs/2211.15455}{arXiv:2211.15455}.
A modern source in the classical arithmetic-undecidability framework.

\bibitem{repo-catalogue}
V. Reshetnikov, \texttt{Surreal} repository, \emph{The surreal and surcomplex
reports}, \texttt{docs/README.md}, inspected 23 September 2026 at commit
\repoSHA.
\url{https://github.com/VladimirReshetnikov/Surreal/blob/958b5c4865819bd55ea1f5ffc050282aea7ef570/docs/README.md}.

\bibitem{repo-quotients}
\texttt{Surreal} repository, \emph{Set-Sized Quotients of the Omnific Integers},
AI-assisted research report, \texttt{article.tex}, same inspected commit.
\url{https://github.com/VladimirReshetnikov/Surreal/blob/958b5c4865819bd55ea1f5ffc050282aea7ef570/docs/surreal/set-sized-quotients-of-omnific-integers/article.tex}.
Cited for related scope and notation, not as a dependency of the new proofs.

\bibitem{repo-exp}
\texttt{Surreal} repository, \emph{Cofinal displacement and the rigidity of
exponential automorphisms}, report README, same inspected commit.
\url{https://github.com/VladimirReshetnikov/Surreal/blob/958b5c4865819bd55ea1f5ffc050282aea7ef570/docs/surreal/exponential-automorphism-rigidity/README.md}.
Cited for related scope and the distinction between generic algebra and concrete
surreal instantiation.
\end{thebibliography}
\end{document}
